\section{Clustering and the EM Algorithm}

\subsection{Likelihood and Log-Likelihood in the Clustering Context}

In the clustering setting, the observations
\[
\underline{X}=(X_1,\ldots,X_n)
\]
are observed, but the group memberships
\[
\underline{Z}=(Z_1,\ldots,Z_n)
\]
are \hl{unobserved (latent)}.

For a Gaussian mixture model,
\[
X_i\mid Z_i=k
\sim
\mathcal{N}(\mu_k,\Sigma_k),
\]
with mixing proportions
\[
P(Z_i=k)=p_k,
\qquad
\sum_{k=1}^Kp_k=1.
\]

The density of an observation $x_i$ is therefore obtained by
summing over the possible groups:
\[
f(x_i;\theta)
=
\sum_{k=1}^K
p_k f_k(x_i;\theta_k).
\]

Hence, assuming independent observations,
the \hl{observed-data likelihood} is
\[
L(\theta;\underline{X})
=
\prod_{i=1}^n
\left(
\sum_{k=1}^K
p_k f_k(x_i;\theta_k)
\right).
\]

The corresponding observed-data log-likelihood is
\[
\ell(\theta;\underline{X})
=
\log L(\theta;\underline{X})
=
\sum_{i=1}^n
\log
\left(
\sum_{k=1}^K
p_k f_k(x_i;\theta_k)
\right).
\]

\paragraph{Why is the clustering likelihood harder?}

When $Z_i$ is known, the likelihood contains products of the form
\[
p_kf_k(x_i;\theta_k),
\]
because we know which group generated $x_i$.

In clustering, however, $Z_i$ is unknown. Therefore, for each observation
$x_i$, we must consider all possible groups:
\[
f(x_i;\theta)
=
\sum_{k=1}^K
p_kf_k(x_i;\theta_k).
\]

The logarithm is then applied to this sum:
\[
\log
\left(
\sum_{k=1}^K
p_kf_k(x_i;\theta_k)
\right).
\]

In general,
\[
\log\left(\sum_k a_k\right)
\neq
\sum_k\log(a_k),
\]
which makes the direct maximisation of the observed likelihood
more difficult.

\paragraph{Important distinction.}

The \hl{observed-data likelihood} uses only the observed variables:
\[
\boxed{
L(\theta;\underline{X})
=
\prod_{i=1}^n
\sum_{k=1}^K
p_kf_k(x_i;\theta_k)
}.
\]

The \hl{complete-data likelihood} assumes that the latent group
memberships $Z$ are also available:
\[
\boxed{
L_c(\theta;\underline{X},Z)
=
\prod_{i=1}^n
\prod_{k=1}^K
\left[
p_kf_k(x_i;\theta_k)
\right]^{\tilde Z_{ik}}
}.
\]

The difference between these two likelihoods is the key reason why
EM is introduced.

\subsection{MLE in the Clustering Context}

The maximum likelihood estimator in clustering is defined by
\[
\boxed{
\hat\theta
=
\arg\max_{\theta}
L(\theta;\underline{X})
}
\]
or, equivalently,
\[
\boxed{
\hat\theta
=
\arg\max_{\theta}
\ell(\theta;\underline{X})
}.
\]

For a Gaussian mixture model,
\[
\hat\theta
=
\arg\max_{\theta}
\sum_{i=1}^n
\log
\left(
\sum_{k=1}^K
p_kf_k(x_i;\mu_k,\Sigma_k)
\right).
\]

The difficulty is that the latent variables $Z_i$ are unknown.

If the memberships were known, we could directly maximise the
complete-data likelihood and obtain the empirical estimators
\[
\hat p_k=\frac{n_k}{n},
\]
\[
\hat\mu_k
=
\frac{1}{n_k}
\sum_{i=1}^n
\tilde Z_{ik}x_i,
\]
and
\[
\hat\Sigma_k
=
\frac{1}{n_k}
\sum_{i=1}^n
\tilde Z_{ik}
(x_i-\hat\mu_k)
(x_i-\hat\mu_k)^T.
\]

However, in clustering the $\tilde Z_{ik}$ are unknown.

Therefore, we cannot directly use these formulas.

\[
\boxed{
\text{The problem is that the quantities needed to estimate the
parameters are themselves unknown.}
}
\]

The Expectation-Maximisation algorithm solves this problem by
alternating between estimating the latent memberships and updating
the model parameters.

\subsection{From the Observed Likelihood to the Completed Likelihood}

The central idea of EM is to work with the
\hl{complete-data likelihood}, because it is easier to maximise.

The complete-data likelihood is
\[
L_c(\theta;\underline{X},Z)
=
P(\underline{X},\underline{Z};\theta).
\]

Using the one-hot representation $\tilde Z_{ik}$,
\[
L_c(\theta;\underline{X},Z)
=
\prod_{i=1}^n
\prod_{k=1}^K
\left[
p_kf_k(x_i;\theta_k)
\right]^{\tilde Z_{ik}}.
\]

Its logarithm is
\[
\ell_c(\theta;\underline{X},Z)
=
\sum_{i=1}^n
\sum_{k=1}^K
\tilde Z_{ik}
\left[
\log p_k+\log f_k(x_i;\theta_k)
\right].
\]

This is much easier to handle because there is no logarithm of a sum.

However, $Z$ is unknown.

Therefore, instead of maximising
\[
\ell_c(\theta;\underline{X},Z)
\]
directly, EM uses its conditional expectation given the observed data.

For a current parameter estimate $\theta^{(t)}$, define
\[
Q(\theta\mid\theta^{(t)})
=
\mathbb{E}_{Z\mid X,\theta^{(t)}}
\left[
\ell_c(\theta;X,Z)
\right].
\]

The algorithm then alternates between:

\[
\boxed{
\text{E-step: estimate the latent variables}
}
\]

and

\[
\boxed{
\text{M-step: maximise the expected complete-data log-likelihood}
}.
\]

\subsection{The EM Algorithm}

Let
\[
\theta^{(t)}
\]
be the current estimate of the model parameters.

\subsubsection{E-Step: Expectation}

The goal is to compute the conditional distribution of the latent
variables given the observations and the current parameters.

For each observation $i$ and group $k$, define
\[
\tau_{ik}^{(t)}
=
P(Z_i=k\mid X_i=x_i,\theta^{(t)}).
\]

Using Bayes' rule,
\[
\tau_{ik}^{(t)}
=
\frac{
p_k^{(t)}
f_k(x_i;\theta_k^{(t)})
}{
\sum_{\ell=1}^K
p_\ell^{(t)}
f_\ell(x_i;\theta_\ell^{(t)})
}.
\]

These quantities are called the
\hl{responsibilities}.

They satisfy
\[
0\leq\tau_{ik}^{(t)}\leq1
\]
and
\[
\sum_{k=1}^K\tau_{ik}^{(t)}=1.
\]

Intuitively,
\[
\boxed{
\tau_{ik}^{(t)}
=
\text{probability that observation $i$ belongs to group $k$}
}
\]
according to the current model.

The E-step can therefore be viewed as replacing the unknown hard
membership indicators $\tilde Z_{ik}$ by their conditional
expectations:
\[
\boxed{
\mathbb{E}[\tilde Z_{ik}\mid X,\theta^{(t)}]
=
\tau_{ik}^{(t)}
}.
\]

Hence,
\[
Q(\theta\mid\theta^{(t)})
=
\sum_{i=1}^n
\sum_{k=1}^K
\tau_{ik}^{(t)}
\left[
\log p_k+\log f_k(x_i;\theta_k)
\right].
\]

\subsubsection{M-Step: Maximisation}

We now maximise the expected complete-data log-likelihood:
\[
\theta^{(t+1)}
=
\arg\max_\theta
Q(\theta\mid\theta^{(t)}).
\]

For a Gaussian mixture, this gives
\[
p_k^{(t+1)}
=
\frac{1}{n}
\sum_{i=1}^n
\tau_{ik}^{(t)},
\]
\[
\mu_k^{(t+1)}
=
\frac{
\sum_{i=1}^n
\tau_{ik}^{(t)}x_i
}{
\sum_{i=1}^n
\tau_{ik}^{(t)}
},
\]
and
\[
\Sigma_k^{(t+1)}
=
\frac{
\sum_{i=1}^n
\tau_{ik}^{(t)}
(x_i-\mu_k^{(t+1)})
(x_i-\mu_k^{(t+1)})^T
}{
\sum_{i=1}^n
\tau_{ik}^{(t)}
}.
\]

Thus, the observations are no longer assigned completely to one
group. Instead, each observation contributes to every group with
weight $\tau_{ik}^{(t)}$.

\subsubsection{Complete EM Iteration}

The EM algorithm is therefore

\[
\boxed{
\theta^{(t)}
\overset{\text{E-step}}{\longrightarrow}
\tau_{ik}^{(t)}
\overset{\text{M-step}}{\longrightarrow}
\theta^{(t+1)}
}.
\]

More explicitly:

\begin{enumerate}
    \item Choose an initial parameter value $\theta^{(0)}$.
    
    \item \textbf{E-step:} compute the responsibilities
    \[
    \tau_{ik}^{(t)}
    =
    P(Z_i=k\mid X_i=x_i,\theta^{(t)}).
    \]
    
    \item \textbf{M-step:} maximise
    \[
    Q(\theta\mid\theta^{(t)})
    \]
    and obtain $\theta^{(t+1)}$.
    
    \item Repeat the E-step and M-step until convergence.
\end{enumerate}

A common stopping criterion is based on the observed log-likelihood:
\[
\left|
\ell(\theta^{(t+1)})
-
\ell(\theta^{(t)})
\right|
<
\varepsilon.
\]

\subsection{Proof of EM Monotonicity and Convergence}

We now prove the fundamental property of EM:
\[
\boxed{
\ell(\theta^{(t+1)})
\geq
\ell(\theta^{(t)})
}.
\]

Thus, the observed log-likelihood does not decrease from one EM
iteration to the next.

\subsubsection*{Step 1: The complete-data distribution}

Let
\[
p_\theta(x,z)
\]
denote the joint density of the observed and latent variables.

Then
\[
p_\theta(x,z)
=
p_\theta(z\mid x)\,p_\theta(x).
\]

Taking logarithms gives
\[
\log p_\theta(x,z)
=
\log p_\theta(z\mid x)
+
\log p_\theta(x).
\]

For fixed observed data $x$, define
\[
Q(\theta\mid\theta^{(t)})
=
\mathbb E_{\theta^{(t)}}
\left[
\log p_\theta(X,Z)
\mid X=x
\right].
\]

Therefore,
\[
Q(\theta\mid\theta^{(t)})
=
\mathbb E_{\theta^{(t)}}
\left[
\log p_\theta(Z\mid X)
\mid X=x
\right]
+
\log p_\theta(x).
\]

Similarly,
\[
Q(\theta^{(t)}\mid\theta^{(t)})
=
\mathbb E_{\theta^{(t)}}
\left[
\log p_{\theta^{(t)}}(Z\mid X)
\mid X=x
\right]
+
\log p_{\theta^{(t)}}(x).
\]

\subsubsection*{Step 2: Subtract the two expressions}

Subtracting gives
\[
Q(\theta\mid\theta^{(t)})
-
Q(\theta^{(t)}\mid\theta^{(t)})
\]
\[
=
\mathbb E_{\theta^{(t)}}
\left[
\log
\frac{
p_\theta(Z\mid X)
}{
p_{\theta^{(t)}}(Z\mid X)
}
\;\middle|\;
X=x
\right]
+
\ell(\theta;x)-\ell(\theta^{(t)};x).
\]

By definition of the Kullback--Leibler divergence,
\[
\operatorname{KL}
\left(
p_{\theta^{(t)}}(Z\mid X)
\;\middle\|\;
p_\theta(Z\mid X)
\right)
\]
\[
=
\mathbb E_{\theta^{(t)}}
\left[
\log
\frac{
p_{\theta^{(t)}}(Z\mid X)
}{
p_\theta(Z\mid X)
}
\;\middle|\;
X=x
\right].
\]

Hence,
\[
\mathbb E_{\theta^{(t)}}
\left[
\log
\frac{
p_\theta(Z\mid X)
}{
p_{\theta^{(t)}}(Z\mid X)
}
\;\middle|\;
X=x
\right]
=
-
\operatorname{KL}
\left(
p_{\theta^{(t)}}(Z\mid X)
\;\middle\|\;
p_\theta(Z\mid X)
\right).
\]

Therefore,
\[
\boxed{
\ell(\theta;x)-\ell(\theta^{(t)};x)
=
Q(\theta\mid\theta^{(t)})
-
Q(\theta^{(t)}\mid\theta^{(t)})
+
\operatorname{KL}
\left(
p_{\theta^{(t)}}(Z\mid X)
\;\middle\|\;
p_\theta(Z\mid X)
\right)
}.
\]

\subsubsection*{Step 3: Use the non-negativity of KL divergence}

Recall that
\[
\operatorname{KL}(P\|Q)\geq0.
\]

Thus,
\[
\ell(\theta;x)-\ell(\theta^{(t)};x)
\geq
Q(\theta\mid\theta^{(t)})
-
Q(\theta^{(t)}\mid\theta^{(t)}).
\]

In the M-step, $\theta^{(t+1)}$ is chosen such that
\[
Q(\theta^{(t+1)}\mid\theta^{(t)})
\geq
Q(\theta^{(t)}\mid\theta^{(t)}).
\]

Therefore,
\[
\ell(\theta^{(t+1)};x)
-
\ell(\theta^{(t)};x)
\geq0.
\]

Hence,
\[
\boxed{
\ell(\theta^{(t+1)};x)
\geq
\ell(\theta^{(t)};x)
}.
\]

This proves that EM produces a
\hl{non-decreasing sequence of observed log-likelihood values}.

\subsubsection*{Step 4: Consequence for the likelihood sequence}

Define
\[
\ell_t
=
\ell(\theta^{(t)};x).
\]

From the previous result,
\[
\ell_{t+1}\geq\ell_t.
\]

Thus,
\[
\ell_0\leq\ell_1\leq\ell_2\leq\cdots.
\]

So the sequence $(\ell_t)$ is monotone non-decreasing.

If the log-likelihood is bounded above, then by the monotone
convergence theorem for real sequences,
\[
\boxed{
\ell_t\longrightarrow\ell^*
}
\]
for some finite value $\ell^*$.

Therefore, under these conditions, the observed-data likelihood
converges.

\subsubsection*{Step 5: Relation with stationary points}

Suppose that the parameter space is such that:

\begin{itemize}
    \item the likelihood is bounded above;
    \item the relevant parameter set is compact, or the sequence
    $\theta^{(t)}$ remains in a compact subset;
    \item the likelihood and $Q$ are sufficiently regular
    (for example, continuous and differentiable where required);
    \item the M-step finds an exact maximum of
    $Q(\theta\mid\theta^{(t)})$.
\end{itemize}

Then the parameter sequence has at least one accumulation point.

Let
\[
\theta^*
\]
be an accumulation point of the EM sequence.

Because the observed log-likelihood is non-decreasing and converges,
the increase in likelihood must vanish asymptotically:
\[
\ell(\theta^{(t+1)})
-
\ell(\theta^{(t)})
\longrightarrow0.
\]

From
\[
\ell(\theta^{(t+1)})-\ell(\theta^{(t)})
\geq
Q(\theta^{(t+1)}\mid\theta^{(t)})
-
Q(\theta^{(t)}\mid\theta^{(t)}),
\]
we also obtain that asymptotically there can be no positive gain in
the $Q$ function.

Therefore, at an accumulation point $\theta^*$, the M-step can no
longer improve $Q$.

Under the regularity assumptions above, this implies
\[
\nabla_\theta Q(\theta^*\mid\theta^*)=0.
\]

The EM fixed-point relation then gives
\[
\nabla_\theta \ell(\theta^*;x)=0.
\]

Hence,
\[
\boxed{
\theta^*
\text{ is a stationary point of the observed-data log-likelihood.}
}
\]

Therefore, the standard convergence statement is:

\[
\fbox{
\parbox{0.8\textwidth}{
Under suitable regularity conditions, every accumulation point
of the EM sequence is a stationary point of the likelihood.}}
\]

\paragraph{Important caveat for Gaussian mixtures.}

For an unrestricted Gaussian mixture model, the observed likelihood is
not necessarily bounded above.

For example, a covariance matrix can collapse around a single
observation:
\[
\Sigma_k\longrightarrow0.
\]

In this situation, the corresponding Gaussian density can become
arbitrarily large, and the likelihood can diverge.

Thus, for an unconstrained GMM, one must be careful with the statement
``EM converges to the MLE.''

The precise statement is that EM increases the likelihood at every
iteration, while convergence to a finite stationary point requires
appropriate regularity conditions or constraints on the parameter
space.

\subsection{Classification EM (CEM)}

The standard EM algorithm keeps the memberships
\[
\tau_{ik}
=
P(Z_i=k\mid X_i=x_i)
\]
as \hl{soft assignments}.

Thus, an observation can belong partially to several groups.

Classification EM (CEM) introduces a
\hl{hard classification step}.

After the E-step, each observation is assigned to the group with the
largest posterior probability:
\[
\hat Z_i
=
\arg\max_{k}
\tau_{ik}.
\]

Equivalently,
\[
\tilde Z_{ik}
=
\begin{cases}
1
&
\text{if }
k=\arg\max_{\ell}\tau_{i\ell},\\
0
&
\text{otherwise}.
\end{cases}
\]

The algorithm therefore becomes

\[
\boxed{
\text{E-step}
\longrightarrow
\text{Classification step}
\longrightarrow
\text{M-step}.
}
\]

In the M-step, the parameters are updated using the hard assignments,
for example
\[
\hat p_k=\frac{n_k}{n},
\]
\[
\hat\mu_k
=
\frac{1}{n_k}
\sum_{i:\hat Z_i=k}x_i.
\]

\paragraph{Properties of CEM.}

CEM has the following characteristics:

\begin{itemize}
    \item It produces a \hl{hard partition} of the observations.
    \item It is usually more directly focused on classification than
    standard EM.
    \item The algorithm can be faster and easier to interpret because
    every observation is assigned to one group.
    \item The objective associated with the hard classification is not
    exactly the same as the ordinary observed-data likelihood of
    standard EM.
    \item Therefore, the monotonicity guarantee of standard EM for the
    observed-data likelihood does not automatically carry over to CEM.
\end{itemize}

Conceptually,
\[
\boxed{
\text{EM: soft membership}
\qquad\qquad
\text{CEM: hard membership}.
}
\]

\subsection{Stochastic EM (SEM)}

Stochastic EM (SEM) replaces the deterministic classification step by
a random sampling step.

After the E-step, the conditional probabilities
\[
\tau_{ik}
=
P(Z_i=k\mid X_i=x_i,\theta^{(t)})
\]
are used to randomly generate a group assignment:
\[
Z_i^{(t)}
\sim
\mathcal{M}
\left(
1,
\tau_{i1}^{(t)},\ldots,\tau_{iK}^{(t)}
\right).
\]

Thus, an observation is assigned to exactly one group, but the
assignment is random.

The algorithm is therefore

\[
\boxed{
\text{E-step}
\longrightarrow
\text{Random sampling of }Z
\longrightarrow
\text{M-step}.
}
\]

\paragraph{Properties of SEM.}

\begin{itemize}
    \item SEM produces \hl{hard classifications}, like CEM.
    \item The classifications are stochastic rather than deterministic.
    \item The random component can help the algorithm explore
    different regions of the parameter space and potentially escape
    some poor local solutions.
    \item Unlike standard EM, the observed log-likelihood is not
    necessarily non-decreasing at every iteration.
    \item The sequence of parameter estimates is therefore stochastic.
\end{itemize}

The main distinction is

\[
\boxed{
\begin{array}{ccc}
\text{EM} &:& \text{soft deterministic memberships}\\
\text{CEM} &:& \text{hard deterministic memberships}\\
\text{SEM} &:& \text{hard stochastic memberships}.
\end{array}
}
\]

\subsection{$k$-Means as a Particular Case of a Gaussian Mixture Model}

The $k$-means algorithm can be interpreted as a particular restricted
case of Gaussian-mixture clustering.

Consider a Gaussian mixture model with

\[
X\mid Z=k
\sim
\mathcal{N}(\mu_k,\sigma^2I),
\]
where

\[
\Sigma_k=\sigma^2I
\qquad
\text{for every }k,
\]
and suppose that the mixing proportions are equal:
\[
p_1=\cdots=p_K=\frac{1}{K}.
\]

Thus, all groups have:

\begin{itemize}
    \item the same spherical covariance matrix;
    \item the same variance in every direction;
    \item the same mixing proportion.
\end{itemize}

For a given observation $x_i$, the posterior probability of group $k$
is proportional to
\[
p_kf_k(x_i).
\]

Because all $p_k$ are equal and all covariance matrices are equal,
maximising the posterior probability is equivalent to maximising
the Gaussian density:
\[
\hat Z_i
=
\arg\max_k f_k(x_i).
\]

For a spherical Gaussian,
\[
f_k(x_i)
\propto
\exp
\left(
-\frac{1}{2\sigma^2}
\|x_i-\mu_k\|^2
\right).
\]

Since the exponential function is increasing,
\[
\hat Z_i
=
\arg\min_k
\|x_i-\mu_k\|^2.
\]

Therefore,

\[
\boxed{
\text{assign each observation to its nearest mean}.
}
\]

This is exactly the assignment step of $k$-means.

\subsubsection{The $k$-Means Objective}

The $k$-means algorithm minimises the within-cluster sum of squared
distances:
\[
\boxed{
J(\mu,Z)
=
\sum_{i=1}^n
\|x_i-\mu_{Z_i}\|^2
}.
\]

The two $k$-means steps are:

\paragraph{Assignment step.}
Given the current means,
\[
Z_i
=
\arg\min_k
\|x_i-\mu_k\|^2.
\]

\paragraph{Update step.}
Given the current assignments,
\[
\mu_k
=
\frac{1}{n_k}
\sum_{i:Z_i=k}x_i.
\]

These are exactly the same two types of operations that appear in
hard Gaussian-mixture estimation:

\[
\boxed{
\begin{array}{ccc}
\text{Assignment} &\longleftrightarrow& \text{Classification}\\
\text{Mean update} &\longleftrightarrow& \text{M-step}.
\end{array}
}
\]

\subsubsection{Connection Between GMM, CEM and $k$-Means}

The relationship can therefore be summarised as follows.

For a general GMM:
\[
X\mid Z=k\sim\mathcal N(\mu_k,\Sigma_k),
\]
the memberships are soft:
\[
\tau_{ik}
=
P(Z_i=k\mid X_i=x_i).
\]

For CEM, these responsibilities are converted into hard assignments:
\[
Z_i
=
\arg\max_k\tau_{ik}.
\]

For $k$-means, we additionally impose spherical and common covariance,
equal mixing proportions, and use the nearest-mean assignment:
\[
Z_i
=
\arg\min_k
\|x_i-\mu_k\|^2.
\]

Thus,
\[
\fbox{
\parbox{0.8\textwidth}{
\(k\)-means can be viewed as a hard, highly constrained
Gaussian-mixture clustering procedure.
}}
\]

More precisely, with
\[
\Sigma_k=\sigma^2I
\]
for all $k$ and equal mixing proportions, the Gaussian posterior
classification rule becomes nearest-centroid classification.

The variance $\sigma^2$ controls the scale of the Gaussian densities.
The $k$-means objective is recovered from the Gaussian log-likelihood
up to multiplicative and additive constants when the common variance
is fixed.

\subsection*{Summary}

The main ideas of this chapter can be summarised by the chain

\[
\boxed{
\begin{array}{c}
\text{Only }X\text{ observed}\\
\Downarrow\\
\text{Latent }Z\text{ variables}\\
\Downarrow\\
\text{Observed likelihood contains a sum over groups}\\
\Downarrow\\
\text{Direct MLE is difficult}\\
\Downarrow\\
\text{Use complete-data likelihood}\\
\Downarrow\\
\text{Take its conditional expectation}\\
\Downarrow\\
\text{EM: E-step + M-step}
\end{array}
}
\]

The three main algorithms are

\[
\boxed{
\begin{array}{lll}
\text{EM} &:& \text{soft memberships}\\
\text{CEM} &:& \text{hard deterministic memberships}\\
\text{SEM} &:& \text{hard stochastic memberships}.
\end{array}
}
\]

And the main connection with $k$-means is

\[
\boxed{
\text{$k$-means}
\subset
\text{hard Gaussian-mixture clustering}
\subset
\text{Gaussian-mixture clustering}.
}
\]

The key distinction is that a general GMM learns both the
cluster means and covariance structures and uses probabilistic
memberships, whereas $k$-means uses nearest-centroid hard assignments
and a common spherical geometry.